%% Thesis language: english or polish
%% Thesis type: master or engineer
\documentclass[polish,master]{polsl-msth}
%% set document encoding
\usepackage[latin2]{inputenc}

%% thesis author
\author{Jan Kowalski}

%% thesis supervisor
\supervisor{Jan Kowalski}

%% thesis consultant, optional
%% \consultant{Jan Kowalski}

%% Thesis title
\title{Przyk�ad wykorzystania klasy polsl-msth}
\begin{document}
\frontmatter
\maketitle
\tableofcontents
\listoftables
\listoffigures
\mainmatter
\chapter{Wprowadzenie}
Dokument ten opisuje wykorzystanie klasy polsl-msth do przygotowania
pracy magisterskiej na Wydziale Automatyki, Elektroniki i~Informatyki
Politechniki �l�skiej. Sam dokument jest
te� zgodny z~szablonem pracy magisterskiej.

W~rozdziale \ref{zalecenia} znajduj� si� zalecenia odno�nie formatowania
pracy.

\chapter{Zalecenia\label{zalecenia}}
Pierwszy akapit rozdzia�u, sekcji lub podsekcji powinien rozpoczyna�
si� bez wci�cia. Wszystkie kolejne powinny mie� wci�cie skonfigurowane
w~klasie.

Opis rysunk�w powinien znajdowa� si� pod rysunkiem jak w~przypadku
rysunku \ref{pi}.
\begin{figure}[hbtp]
\centerline{\includegraphics[width=1in]{pi}}
\caption{Logo Politechniki �l�skiej.\label{pi}}
\end{figure}

Opis tabel powinien znajdowa� si� nad tabel� jak w~przypadku
tabeli \ref{tab}.
\begin{table}[hbtp]
\caption{Przyk�adowa tabela.\label{tab}}
\centerline{\begin{tabular}{|c|c|}
\hline
1 & 2\\
\hline
3 & 4\\
\hline
\end{tabular}}
\end{table}

R�wnania powinny by� numerowane jak r�wnanie \ref{entropy}.
\begin{equation}
H(X) = -\sum_{i=0}^n P(x_i) \log_2 P(x_i).\label{entropy}
\end{equation}
\end{document}
